% !PRINTED 309
\weberpage{309}
%<<CONT
so erhalten die Functionen (1), von positiven Zahlenfactoren abgesehen, für $x = 0$ die Werthe
\weq{2}{a_n, a_{n-1}, a_{n-2} \ldots a_0,}
und daraus ergiebt sich also das Theorem:

\gesperrt{IV. Die Anzahl der positiven Wurzeln der Gleichung $f(x) = 0$ ist gleich oder um eine gerade Zahl kleiner, als die Anzahl der Zeichenwechsel in der Reihe der Coëfficienten von $f(x)$, wobei etwa verschwindende Coëfficienten einfach zu übergehen sind und die Wurzel $x = 0$, wenn sie vorhanden ist, nicht mitzählt.}

Um auch die Anzahl der negativen Wurzeln abzuschätzen, kann man entweder das \gesperrt{Budan-Fourier}'sche Theorem auf das Intervall $(-\infty, 0)$ anwenden, oder man ersetzt $x$ durch $-x$, d. h. man ändert die Vorzeichen von $a_1, a_3, a_5 \ldots$ und wendet dann das Theorem IV. an.

In ähnlicher Weise wollen wir das \gesperrt{Newton}'sche Kriterium des \S.\,100 specialisiren.

Wir bemerken dazu, dass
\[
f_\nu(0) = \frac{\Pi(n-\nu)\,\Pi(\nu)}{\Pi(n)}\, a_{n-\nu},
\]
und folglich
\[
F_\nu(0) = \frac{\Pi(\nu)\,\Pi(\nu-1)\,\Pi(n-\nu)\,\Pi(n-\nu-1)}{\Pi(n)\,\Pi(n)}\left[\nu(n-\nu)\,a_{n-\nu}^2 - (\nu+1)(n-\nu+1)\,a_{n-\nu-1}\,a_{n-\nu+1}\right]
\]
wird, dass ferner die Reihe der $f_\nu(x)$ für $x = +\infty$ nur Zeichenfolgen, für $x = -\infty$ nur Zeichenwechsel darbietet.

Wenn wir also
\weq{3}{A_\nu = \nu(n-\nu)\,a_\nu^2 - (\nu+1)(n-\nu+1)\,a_{\nu-1}\,a_{\nu+1},}
\[
A_0 = 1, \quad A_n = 1,
\]
setzen, so dass $A_\nu$, von einem positiven Factor abgesehen, mit $F_{n-\nu}(0)$ übereinstimmt, so folgt, dass für $x = 0$ die Doppelreihe (7), \S.\,100, in umgekehrter Ordnung geschrieben, den Vorzeichen nach übereinstimmt mit
\weq{4}{
\begin{array}{c}
a_0,\ a_1,\ a_2 \ldots a_n \\
A_0,\ A_1,\ A_2 \ldots A_n.
\end{array}
}
Beim Uebergang zu $x = +\infty$ sind in der Reihe $f(x), f_1(x) \ldots f_n(x)$ alle Zeichenwechsel und beim Uebergang zu
%CONT>>
